\documentclass[11pt]{article}
\usepackage{amsmath,amssymb}
\usepackage[margin=1in]{geometry}
\title{Room of Mathematicians, item 6 (P6 $\beta_2$ irrationality) --- Seat: Erd\H{o}s}
\author{}
\date{2026-09-26}
\begin{document}
\maketitle

\noindent\textbf{Task.} Widen the PSLQ/LLL integer-relation search for $q^\ast$, $\beta_2=1/\sqrt{q^\ast}$,
$\tau^\ast=-\ln(q^\ast)/\pi$ beyond the basis already excluded
($\{\pi,e,\sqrt2,\sqrt3,\sqrt5,\varphi,\log2,\log3,\gamma,\text{Catalan},\zeta(3)\}$, and beyond
the degree-$\le128$ LLL polynomial-relation exclusion, both of which are NOT re-run here per
instructions).

\section{Calibration point (reproduced first)}

Read \texttt{code/zeta\_probe/rootunity\_zeros/general/P6/P6\_note.tex} and
\texttt{P6\_v2/P6\_v2\_note.tex} in full. Using the pre-existing 2090-digit certified bracket at
\texttt{private/ROOM/Beta2/g1\_certificate/bracket.txt} (lines 1--2 give $q^\ast$'s digit string to
2090 digits, lines 4--5 give $\beta_2$'s to 2080 digits), reproduced at 120-digit working precision
in this session (mpmath, \texttt{mp.dps=130}):
\[
q^\ast = 0.449453630558948046125545825395706089225112808545993877680332962338142926481334934297071480047\ldots
\]
\[
\beta_2 = 1.491617787114374226836712746983063199653295791260103719567394221570556692679336\ldots
\]
Cross-check: $\beta_2 - 1/\sqrt{q^\ast} \approx -4.7\times10^{-80}$ at 130-digit working
precision --- consistent, digit-for-digit, with the bracket file and with memory note
\texttt{route-b-singularity}. Confirmed (matches P6\_v2\_note.tex \S2): \texttt{mpmath.identify}
already found no match for $q^\ast$, $\beta_2$, or $\tau^\ast=-\ln(q^\ast)/\pi$ against
$\{\pi,e,\sqrt2,\sqrt3,\sqrt5,\varphi,\log2,\log3,\gamma,\text{Catalan},\zeta(3)\}$. This exact
basis, and the degree-$\le64/128$ LLL exclusion, were \textbf{not} retried here per instructions.

\section{New, wider bases tried (PSLQ, mpmath, 100-digit working precision, maxcoeff $10^8$,
maxsteps 5000, each call under the \texttt{perl alarm 290} wrapper)}

\begin{enumerate}
\item \textbf{Gamma special values.} $\{1,\ \Gamma(1/3),\ \Gamma(1/4),\ \Gamma(1/6),\
\Gamma(1/4)^2/\sqrt\pi\ (\text{lemniscate constant}),\ \Gamma(1/3)^3\}$ against $q^\ast$,
$\beta_2$, and $\tau^\ast$ separately (3 PSLQ calls). \textbf{No relation found} (all returned
\texttt{None}).
\item \textbf{$q$-Pochhammer/eta-adjacent constants at $q=q^\ast$ itself.} Computed
$(q^\ast;q^\ast)_\infty=\prod_{n\ge1}(1-q^{\ast n})$ directly (400-term product, converged to
working precision) and the formal eta-like combination $q^{\ast-1/24}(q^\ast;q^\ast)_\infty$
(this is only a formal analogue --- $q^\ast$ is a real base variable here, not a nome
$e^{2\pi i\tau}$ for algebraic $\tau$; P6\_v2\_note.tex \S1--2 already showed $q^\ast$ is not of
CM-nome form). Basis $\{1,\ (q^\ast;q^\ast)_\infty,\ \text{eta-like}\}$ against $q^\ast$ and
$\beta_2$ (2 calls). \textbf{No relation found.}
\item \textbf{Bessel zero constants.} $\{1,\ j_{0,1},\ j_{1,1}\}$ (first zeros of ordinary Bessel
$J_0,J_1$, via \texttt{mpmath.besseljzero}) against $q^\ast$ and $\beta_2$ (2 calls). This was
flagged as a long shot given the Hahn--Exton $q$-Bessel origin of $S(q)$; the classical-Bessel
zeros are a different (non-$q$-deformed) object and there was no strong a priori reason to expect
a hit. \textbf{No relation found.}
\item \textbf{Algebraic self-combinations of $\beta_2$.} $\{1,\ \beta_2,\ \sqrt{\beta_2},\
\beta_2^{1/3},\ 1/\beta_2,\ \beta_2-1,\ \pi,\ e,\ \sqrt2\}$ in one 9-term PSLQ call. PSLQ
\emph{did} return a nonzero relation: coefficients $[1,-1,0,0,0,1,0,0,0]$, i.e.
$1 - \beta_2 + (\beta_2-1) = 0$. \textbf{This is a spurious/trivial hit}: it is the tautology
$1-\beta_2+(\beta_2-1)\equiv0$ coming from including $\beta_2-1$ as its own basis vector, not a
relation involving $\pi,e,\sqrt2$ or any new constant. Flagged explicitly so it is not
mis-read as a discovery; PSLQ correctly finds the shortest relation in its search space, and
that happens to be the planted algebraic identity among the self-combination terms.
\item \textbf{Wide direct re-check.} $\{1,\pi,e,\sqrt2,\sqrt3,\sqrt5\}$ against $q^\ast$ directly,
and $\{1,\pi,e,\sqrt2,\sqrt3,\sqrt5,\log2\}$ against $\tau^\ast$ (2 calls, slightly different
vector packing than the pre-existing \texttt{identify} check, included as a second-opinion
cross-check with the raw PSLQ call rather than \texttt{mpmath.identify}'s wrapper).
\textbf{No relation found}, consistent with the pre-existing result.
\end{enumerate}

\section{Verdict}

\textbf{NOTHING FOUND.} All genuine (non-tautological) PSLQ searches against the four wider,
more exotic bases --- Gamma special values ($\Gamma(1/3),\Gamma(1/4),\Gamma(1/6)$, lemniscate
constant), $q$-Pochhammer/eta-adjacent constants formed from $q^\ast$ itself, ordinary-Bessel
first zeros, and small algebraic self-combinations of $\beta_2$ --- returned no integer relation
at 100-digit working precision, maxcoeff $10^8$. The one nonzero PSLQ output obtained (item 4)
is a tautological identity among the self-combination terms themselves, not a relation to
$\pi,e,\sqrt2$, and is explicitly \emph{not} being reported as a finding.

This is consistent with, and does not go beyond, the existing negative results in
P6\_note.tex/P6\_v2\_note.tex: $q^\ast$/$\beta_2$/$\tau^\ast$ continue to show no recognizable
closed form against every basis tried so far, exotic or standard. As in the prior notes, absence
of a PSLQ hit at 100--130 digits against a finite list of bases is evidence, not proof, of
transcendence or of non-membership in the field generated by these constants; a genuinely new
arithmetic input (per P6\_note.tex \S5) is still what P6 needs, not a wider numerical search.

\section{Caveats}
\begin{itemize}
\item Working precision here (100--130 digits) is far below the pre-existing 2090-digit
certificate; a coincidental near-match at this precision would need re-checking at higher
precision before being taken seriously, but none was found even at this weaker level.
\item The $q$-Pochhammer/eta-like construction (item 2) is only a formal analogue: $q^\ast$ is
not of nome form, so there is no strong theoretical reason to expect $(q^\ast;q^\ast)_\infty$ to
be arithmetically special; the check was included per the task's suggestion, at low cost, not
because of a specific expectation.
\item Bessel-zero basis (item 3) was explicitly a long shot per the task instructions; its
negative result narrows, but does not close, that direction.
\end{itemize}

\end{document}
